\chapter*{Abstract}
\addcontentsline{toc}{chapter}{Abstract}

Efficient fleet management and optimal vehicle dispatch represent critical challenges in modern transportation and logistics systems. Traditional manual dispatch methods are time-consuming, error-prone, and fail to consider multiple optimization factors such as vehicle proximity, historical performance, workload distribution, and operator expertise. This project presents a comprehensive \textbf{Fleet Management System with an Intelligent AI-Powered Dispatch Engine} that addresses these limitations through automated, data-driven vehicle assignment mechanisms.

The system implements a distributed microservices architecture comprising four primary components: a React-based frontend application for real-time visualization and user interaction, a Node.js/Express backend service managing core business logic and API orchestration, a Python/FastAPI machine learning microservice providing intelligent dispatch predictions, and a vehicle simulator for hardware device emulation and system testing. The architecture enables independent scaling, technology diversity, and fault isolation while maintaining seamless inter-service communication through REST APIs, WebSocket connections, and MQTT messaging protocols.

The dispatch engine employs a dual-strategy approach, integrating both rule-based and machine learning-based vehicle selection algorithms. The rule-based engine utilizes a multi-factor weighted scoring algorithm that evaluates five critical parameters: historical performance (10\% weight), current fatigue levels (15\% weight), fault type expertise (25\% weight), criticality matching (5\% weight) and distance (45\% weight). The machine learning component leverages a RandomForest regression model trained on six engineered features—distance metrics, past performance scores, fault history patterns, fatigue calculations, and severity mappings—to predict optimal vehicle-fault assignments with enhanced accuracy.

Real-time GPS tracking and route optimization are facilitated through integration with the Open Source Routing Machine (OSRM) service, implementing a circuit breaker pattern with automatic fallback to Haversine distance calculations for resilience. The system maintains continuous location updates through MQTT-based communication with IoT hardware devices (ESP32/Arduino), enabling bidirectional messaging for dispatch alerts, driver confirmations, and fault resolution notifications. WebSocket technology ensures instantaneous frontend updates for all vehicle movements, status changes, and dispatch events, creating a responsive user experience with minimal latency.

The system implements comprehensive automated workflows spanning the entire fault lifecycle: from initial fault reporting through intelligent vehicle assignment, route calculation, driver confirmation, GPS-guided navigation, arrival detection, and automatic resolution tracking. Role-based access control with JWT authentication ensures secure system access across three permission levels (Administrator, Dispatcher, Viewer), while MongoDB serves as the persistent data store with optimized indexing strategies for high-performance query execution.

Key contributions of this project include: (1) a hybrid dispatch strategy combining deterministic rule-based logic with adaptive machine learning predictions, (2) a resilient microservices architecture with automatic fallback mechanisms ensuring system availability, (3) seamless integration of IoT devices through standardized MQTT messaging protocols, (4) real-time data synchronization across distributed components using event-driven architecture patterns, and (5) comprehensive automated workflow management reducing manual intervention and human error.

The system has been evaluated through prototype testing with simulated vehicle fleets, demonstrating significant improvements in dispatch decision accuracy, response time reduction, and workload distribution optimization compared to traditional manual methods. The implementation provides a scalable, maintainable, and extensible foundation for fleet management operations, with potential applications in emergency services, utility maintenance, logistics coordination, and field service management domains.

\textbf{Keywords:} Fleet Management, Intelligent Dispatch, Machine Learning, Microservices Architecture
